% La derivada como aproximación lineal: jacobiano, gradiente y hessiana (partición por clase: apunte/scripts/reorganizar_por_clase.py)
\chapter{La derivada como aproximación lineal: jacobiano, gradiente y hessiana}\label{cap:derivada}

\section{La idea}


Pasamos del plano a la superficie curva: plano tangente, base local y métrica. Tenemos como horizonte el elipsoide.

\subsection{Pasar del plano a una superficie}\label{sec:plano-superficie}

{\small\vspace{1mm}\linkUnidad{10_jacobiano_primera_forma}{Simulación interactiva: Unidad 14 — Jacobiano y primera forma fundamental}\par}


Sobre una superficie curva no existe una única base cartesiana válida en todas partes. En
cada punto aparece un plano tangente y una base local que cambia al desplazarse. La métrica
indica cómo convertir pequeños cambios de coordenadas en distancias y ángulos reales sobre
la superficie.

\paragraph{Primera traducción algebraica.}
\[
  ds^2=E\,du^2+2F\,du\,dv+G\,dv^2.
\]
Los coeficientes \(E,F,G\) describen cuánto miden y qué ángulo forman los dos vectores de la
base local. La expresión transforma pequeños cambios \((du,dv)\) en una distancia geométrica.

\begin{idea}[De lo global a lo local]
Una superficie curva se estudia mediante pequeños parches que se parecen a un plano. En cada
parche se usan coordenadas locales; al conectar los parches se recupera la geometría global.
Esta es la intuición básica de una variedad.
\end{idea}

\begin{idea}[La derivada como mejor aproximación lineal]
La derivada es la mejor aproximación lineal a una función en un punto dado; el
gradiente es perpendicular a una superficie o curva de nivel porque es
una forma lineal, y la curva de nivel 0 de esa forma lineal es el plano tangente.
\end{idea}

\section{Construcciones, cuentas y ejemplos}


Ponemos el aparato en cuentas: parametrizaciones, espacio tangente, primera forma fundamental, esfera y elipsoide, marcos de referencia y remuestreo. La cátedra aporta los resúmenes y el apunte del elipsoide; las cuentas son desarrollo nuestro.

\subsection{Coordenadas curvilíneas y bases locales}\label{sec:curvilineas}

Una parametrización
\[
  r:U\subset\R^2\longrightarrow\R^3,
  \qquad
  (u,v)\longmapsto r(u,v)
\]
describe una superficie mediante dos coordenadas. Los vectores
\[
  r_u=\frac{\partial r}{\partial u},
  \qquad
  r_v=\frac{\partial r}{\partial v}
\]
forman una base tangente local cuando son linealmente independientes.

\begin{center}
\begin{tikzpicture}[scale=.9]
  \draw[teallight,fill=teallight!60] (1,1)
    .. controls (3,.2) and (6,.5) .. (8,1.4)
    .. controls (7,3.8) and (3.2,4.5) .. (1,3)
    .. controls (.7,2.2) and (.7,1.6) .. cycle;
  \coordinate (p) at (4.6,2.1);
  \fill[ink] (p) circle (2.5pt) node[below left] {$p$};
  \draw[mainline,-{Latex[length=2.6mm]}] (p)--++(2.2,.25)
    node[right] {$r_u$};
  \draw[rust,line width=1.05pt,-{Latex[length=2.6mm]}] (p)--++(-.6,1.7)
    node[above] {$r_v$};
  \draw[muted,line width=.8pt,-{Latex[length=2.6mm]}] (p)--++(.2,2.2)
    node[above] {$n$};
  \node[text=muted] at (6.8,.8) {superficie};
\end{tikzpicture}
\end{center}

En coordenadas cartesianas la base es fija. En una superficie curva, la base cambia de
orientación y de longitud al cambiar el punto. Esta dependencia espacial es la razón por la
cual los productos internos y las derivadas adquieren coeficientes adicionales.

\subsection{Variedades y espacio tangente}\label{sec:variedades}

Una variedad es un objeto que localmente se parece a un espacio euclídeo. Una superficie
regular \(S\subset\R^3\) tiene en cada punto \(p\) un espacio tangente
\[
  T_pS=\operatorname{span}\{r_u(p),r_v(p)\}.
\]
El plano tangente es la mejor aproximación lineal local a la superficie. No es una proyección
cartográfica: pertenece al punto \(p\) y cambia al desplazarse sobre \(S\).

Si \(\Phi:S\to\mathcal N\) es una aplicación entre superficies (variedades), su derivada
\[
  d\Phi_p:T_pS\longrightarrow T_{\Phi(p)}\mathcal N
\]
es una transformación lineal entre espacios tangentes. Esa transformación contiene la
información local sobre escala, orientación, deformación angular y deformación de áreas.

\subsection{Métrica y primera forma fundamental}\label{sec:metrica}

{\small\vspace{1mm}\linkUnidad{10_jacobiano_primera_forma}{Simulación interactiva: Unidad 14 — Jacobiano y primera forma fundamental}\par}


La métrica inducida por la parametrización queda determinada por
\[
  E=\langle r_u,r_u\rangle,\qquad
  F=\langle r_u,r_v\rangle,\qquad
  G=\langle r_v,r_v\rangle.
\]
El elemento de longitud es
\[
  \boxed{ds^2=E\,du^2+2F\,du\,dv+G\,dv^2}.
\]
En forma matricial:
\[
  ds^2=
  \begin{pmatrix}du&dv\end{pmatrix}
  \begin{pmatrix}E&F\\F&G\end{pmatrix}
  \begin{pmatrix}du\\dv\end{pmatrix}.
\]

\paragraph{Esfera.}
Para una esfera de radio \(R\), usando longitud \(\lambda\) y latitud \(\varphi\),
\[
  ds^2=R^2\cos^2\varphi\,d\lambda^2+R^2\,d\varphi^2.
\]
El factor \(\cos\varphi\) explica por qué un incremento fijo de longitud representa una
distancia cada vez menor al acercarse a los polos.

\section{En profundidad}
